Based on the chalkboard notes and the context of your work with Anatole von Lilienfeld's group (U of T / Vector Institute), this makes perfect sense. You are looking at utilizing Euler’s Homogeneous Function Theorem to constrain or simplify the machine learning of interatomic potentials.

Here is a breakdown of the board, the theoretical connection to your paper, and the Python strategy to generate the training data you need.

1. Deciphering Anatole’s Chalkboard

The board is essentially outlining a strategy to link Energies (global quantities) to Forces (gradients) via scaling laws.

Top Right (The Core Theorem):

The equation f(x)=x⋅∂ 
x
​	
 f (with a scaling factor) is Euler's Theorem for Homogeneous Functions.

It states that if a function f satisfies f(s⋅x)=s 
d
 f(x) (it is homogeneous of degree d), then:

x⋅∇f(x)=d⋅f(x)
The "ML linear" arrow: This is the key insight. If your potential is homogeneous, the relationship between the Energy (f) and the Force (∂ 
x
​	
 f) is linear. You don't need to learn them separately; measuring the force gives you the energy (scaled by d).

Left Box (Potential Types):

Coulomb: 1/r (Homogeneous, degree d=−1).

Dispersion: 1/r 
6
  (Homogeneous, degree d=−6, noted on the board as 1/(sx) 
6
 ).

Morse: Listed as a target, but problematic because it contains an exponential, which is not homogeneous in distance r.

Bottom Right (Thermodynamics):

The board draws an analogy to thermodynamics ( 
U
˙
 (S,V,N)...). Extensive thermodynamic potentials are homogeneous first-order functions of extensive variables. Anatole is likely suggesting that atomic potentials can be treated with the same "scaling" logic as thermodynamic equations of state.

2. The Challenge: "Adaptive" Morse Potentials

The Morse potential is defined as:

V(r)=D 
e
​	
 (1−e 
−a(r−r 
e
​	
 )
 ) 
2
 
Unlike Coulomb or Dispersion, this is not homogeneous in r (you cannot pull out a factor s 
d
 ).

The Goal: To apply the "Homogeneous ML" approach to Morse, you likely need to map it to a space where it behaves homogeneously, or learn the "Local Homogeneity" (effective degree d).

Adaptive Strategy: Instead of learning V(r) directly, you might be training a model to predict the effective scaling factor (or "local virial") d(r) where:

d(r)= 
V(r)
r⋅V 
′
 (r)
​	
 
Or, you are using the Law of Corresponding States: by scaling r by r 
e
​	
  and V by D 
e
​	
 , all Morse potentials with the same "stiffness" (ar 
e
​	
 ) collapse to a single curve.

3. Generating the Training Data

To demonstrate this adaptivity, you shouldn't just generate random Morse curves. You need to generate data that scans the parameter space (D 
e
​	
 ,r 
e
​	
 ,a) and calculates the homogeneity violation (or effective degree) to see how well the linear ML model adapts.

Here is a Python script to generate the specific "Training Information" you need. It calculates the Energy, Force, and the "Euler Discrepancy" (how far it is from being homogeneous).

Python
import numpy as np
import pandas as pd
import matplotlib.pyplot as plt

def morse_potential(r, De, a, re):
    """
    Standard Morse potential form.
    """
    return De * (1 - np.exp(-a * (r - re)))**2

def morse_force(r, De, a, re):
    """
    Analytical force (negative gradient) of Morse potential.
    F = -dV/dr
    """
    term = np.exp(-a * (r - re))
    dV_dr = 2 * De * (1 - term) * (term * a)
    return -dV_dr

def generate_adaptive_training_set(n_samples=1000, r_range=(0.8, 5.0)):
    """
    Generates training data for adaptive ML of Morse potentials.
    Focuses on capturing the relationship between V and F via Euler's theorem.
    """
    data = []

    # 1. Sample the Parameter Space (Physics-informed ranges)
    # These ranges cover typical chemical bonds
    des = np.random.uniform(0.5, 5.0, n_samples)   # Well depth (eV)
    res = np.random.uniform(1.0, 3.0, n_samples)   # Equilibrium bond length (Angstrom)
    # 'a' controls the width. Linked to force constant k ~ 2*De*a^2
    as_ = np.random.uniform(0.5, 2.5, n_samples)   

    for i in range(n_samples):
        De, re, a = des[i], res[i], as_[i]
        
        # 2. Adaptive Sampling in r
        # Instead of uniform r, we sample relative to re to capture the well shape
        # This is critical for "Corresponding States" learning
        r_star = np.linspace(0.5, 3.0, 50) # Reduced distance r/re
        rs = r_star * re
        
        Vs = morse_potential(rs, De, a, re)
        Fs = morse_force(rs, De, a, re) # Force = -Gradient
        Gradients = -Fs 

        # 3. Calculate "Local Homogeneity" (The 'Adaptive' Feature)
        # If the function were homogeneous of degree k, then: r * Gradient = k * V
        # We calculate k_eff(r) = (r * Gradient) / V
        # This 'k_eff' is likely what the ML needs to learn to adapt to non-homogeneous potentials.
        
        # Handle division by zero at V=0 (at r=re)
        with np.errstate(divide='ignore', invalid='ignore'):
            k_eff = (rs * Gradients) / Vs
            k_eff[Vs == 0] = 0 # Asymptotically approaches infinity or 0 depending on limit
            
        for r, V, F, g, k in zip(rs, Vs, Fs, Gradients, k_eff):
            data.append({
                'id': i,
                'De': De,
                're': re,
                'a': a,
                'r': r,
                'V': V,          # Target 1: Energy
                'F': F,          # Target 2: Force
                'Euler_LHS': r * g,      # r * dV/dr
                'k_eff': k,      # The effective homogeneity degree at this point
                'stiffness': a * re # Dimensionless shape parameter
            })

    return pd.DataFrame(data)

# --- Execution ---
df = generate_adaptive_training_set(n_samples=50) # 50 potentials * 50 points = 2500 datapoints

# Visualization to check the "Homogeneity Space"
plt.figure(figsize=(10, 6))
# Filter out extreme asymptotes near r_e for plotting clarity
plot_data = df[(df['r'] > df['re']*1.1) & (df['r'] < df['re']*2.5)]

plt.scatter(plot_data['r']/plot_data['re'], plot_data['k_eff'], c=plot_data['stiffness'], cmap='viridis', alpha=0.5)
plt.colorbar(label='Stiffness (a * re)')
plt.title("Effective Homogeneity Degree (k_eff) vs Reduced Distance")
plt.xlabel("Reduced Distance (r / re)")
plt.ylabel("Local Homogeneity Degree (r * V' / V)")
plt.axhline(-6, color='r', linestyle='--', label='Dispersion limit (k=-6)')
plt.axhline(-1, color='b', linestyle='--', label='Coulomb limit (k=-1)')
plt.legend()
plt.grid(True)
plt.show()

print(df.head())
How this helps your paper:

The "Euler_LHS" Column: This is your x * grad_f. By plotting this against V, you can show Anatole exactly where Morse deviates from the linear ML relation he drew.

The "Stiffness" Parameter (a * re): The visualization will likely show that for a fixed stiffness, the "Homogeneity Degree" curve is invariant. This proves that you can adaptively learn the entire Morse family by just learning the mapping of Stiffness → Homogeneity Curve.

Adaptive Learning: You can propose that the ML model predicts k_eff(r) rather than V(r). Since k_eff is smoother and bounded (unlike V which explodes), it is easier to learn. Then you reconstruct V by integrating:

V(r)=V(r 
0
​	
 )exp(∫ 
r 
0
​	
 
r
​	
  
x
k 
eff
​	
 (x)
​	
 dx)
This effectively turns the non-homogeneous Morse problem into a homogeneous-compatible one, aligning with the "linear ML" note on the board.
